\RequirePackage{iftex}
\ifLuaTeX\else
  \GenericError{}{This document must be compiled with LuaLaTeX.}{}{Run `latexmk -lualatex sample.tex` or set your editor's TeX engine to LuaLaTeX.}
  \expandafter\endinput
\fi
\documentclass[9pt]{beamer}
\usepackage{luatexja}
\usepackage{mathtools}
\usepackage{amssymb}
\usepackage{amsthm}
\usepackage{braket}
\usepackage{bm}
\usepackage{booktabs}
\usepackage{array}
\usepackage{tikz}
\usetikzlibrary{arrows.meta, calc}

\setbeamertemplate{proof begin}{}
\setbeamertemplate{proof end}{}

\usetheme[block=fill]{metropolis}
\usefonttheme{professionalfonts}

\setbeamersize{
  text margin left=10mm,
  text margin right=10mm
}

\let\Re\relax
\DeclareMathOperator{\Re}{Re}
\let\Im\relax
\DeclareMathOperator{\Im}{Im}

%%%%%%%%%%%%%%%%%%%%%%%%%%%%%%%% １ページ目 %%%%%%%%%%%%%%%%%%%%%%%%%%%%%%%%%
\title{数学講究前期第9回}
\author{今泉 力}
\date{2026年6月17日}

\begin{document}
\begin{frame}
  \titlepage
\end{frame}

%%%%%%%%%%%%%%%%%%%%%%%%%%%%%%%% ２ページ目 %%%%%%%%%%%%%%%%%%%%%%%%%%%%%%%%%
\begin{frame}
    \frametitle{行列の固有値推定（アダマールテスト）}
    \textbf{アダマールテスト} $\cdots$ ユニタリ行列Uの固有値$e^{i\lambda}$を推定するサブルーチン \\
    第１量子ビットは$\ket{0}$として初期化し, 第２量子ビット以降に状態$\ket{\psi}$を入力する. まず, $\ket{\psi}$が$U$の固有ベクトルである場合を考える. 初期状態は
    \begin{align}
        \ket{0} \otimes \ket{\psi}
    \end{align}
    である.第１量子ビットにアダマールゲートをかけ, 重ね合わせの状態を作る.
    \begin{align}
        \frac{1}{\sqrt{2}} \left( \ket{0} + \ket{1} \right) \otimes \ket{\psi}
    \end{align}
    次に第１量子ビットをコントロールビットとして, 第２量子ビット以降にユニタリ演算$U$を作用させるCtrl-$U$を作用させる.
    \begin{align}
        \begin{bmatrix}
            1 & 0                 &     \\
            0 & 1                 &     \\
              &                   & \text{\huge{$U$}}
        \end{bmatrix}
        \frac{1}{\sqrt{2}} \left( \ket{0} + \ket{1} \right) \otimes \ket{\psi}
        &=
        \frac{1}{\sqrt{2}} \left( \ket{0} \otimes \ket{\psi} + \ket{1} \otimes U\ket{\psi} \right) \\
        &= 
        \frac{1}{\sqrt{2}} \left( \ket{0} \otimes \ket{\psi} + \ket{1} \otimes e^{i\lambda} \ket{\psi} \right) \\
        &= 
        \frac{1}{\sqrt{2}} \left( \ket{0} + e^{i\lambda} \ket{1} \right) \otimes \ket{\psi}
    \end{align}
\end{frame}

%%%%%%%%%%%%%%%%%%%%%%%%%%%%%%%%%%%%%%% ３ページ目 %%%%%%%%%%%%%%%%%%%%%%%%%%%%%%%%%%%%%
\begin{frame}
    \frametitle{行列の固有値推定（アダマールテスト）}
    その後, 第１量子ビットにアダマールゲートをかけると
    \begin{align}
        H\left( \frac{1}{\sqrt{2}} \left( \ket{0} + e^{i\lambda} \ket{1} \right) \otimes \ket{\psi} \right) 
        &= 
        \left( \dfrac{1+e^{i\lambda}}{2} \ket{0} + \dfrac{1-e^{i\lambda}}{2} \ket{1} \right) \otimes \ket{\psi}
    \end{align}
    測定結果が０になる確率$p_0$と１になる確率$p_1$を求めると, ボルンの確率規則により
    \begin{align}
        p_0 = \left| \dfrac{1+e^{i\lambda}}{2} \right|^2 = \dfrac{1+\cos{\lambda}}{2} ,\quad p_1 = \left| \dfrac{1-e^{i\lambda}}{2} \right|^2 = \dfrac{1-\cos{\lambda}}{2}
    \end{align}
    この手法により, どれだけ大きいユニタリ行列$U$であっても, $\cos^-1$を使って$\lambda$の値をもとめることができる.
\end{frame}

%%%%%%%%%%%%%%%%%%%%%%%%%%%%%%%%%%%%%%% ３ページ目 %%%%%%%%%%%%%%%%%%%%%%%%%%%%%%%%%%%%%
\begin{frame}
    \frametitle{行列の固有値推定（アダマールテスト）}
    $\ket{\psi}$が$U$の固有ベクトルでない場合も考える. アダマールゲートを通した後は
    \begin{align}
        \frac{1}{\sqrt{2}} \left( \ket{0} + \ket{1} \right) \otimes \ket{\psi}
    \end{align}
    であり, Ctrl-$U$を作用させると,
    \begin{align}
        \begin{bmatrix}
            1 & 0                 &     \\
            0 & 1                 &     \\
              &                   & \text{\huge{$U$}}
        \end{bmatrix}
        \frac{1}{\sqrt{2}} \left( \ket{0} + \ket{1} \right) \otimes \ket{\psi}
    \end{align}
    最後にアダマールゲート$H$を作用させ,
    \begin{align}
        H \left( \frac{1}{\sqrt{2}} \left( \ket{0} + \ket{1} \right) \otimes \ket{\psi} \right) =
         \ket{0} \otimes \left( \dfrac{\ket{\psi} + U\ket{\psi}}{2} \right) + \ket{1} \otimes \left( \dfrac{\ket{\psi} - U\ket{\psi}}{2} \right)
    \end{align}
\end{frame}

%%%%%%%%%%%%%%%%%%%%%%%%%%%%%%４ぺーじ目%%%%%%%%%%%%%%%%%%%%%%%%%%%%%%%
\begin{frame}
    \frametitle{行列の固有値推定（アダマールテスト）}
    この状態について$p_0$を求めると
    \begin{align}
        p_0 &= \left( \dfrac{\ket{\psi} + U\ket{\psi}}{2} \right)^* \left( \dfrac{\ket{\psi} + U\ket{\psi}}{2} \right) 
        = 
        \frac{1}{4} (\bra{\psi} + \bra{\psi}U^*)(\ket{\psi} + U\ket{\psi}) \\
        &=
        \frac{1}{4} (\braket{\psi}{\psi} + \bra{\psi}U\ket{\psi} + \bra{\psi}U^*\ket{\psi} + \bra{\psi}U^*U\ket{\psi}) \\
        &=
        \frac{1}{4} \left( 2 + \bra{\psi}U\ket{\psi} + (\bra{\psi}U\ket{\psi})^* \right) 
        = 
        \dfrac{1 + \Re \bra{\psi}U\ket{\psi}}{2}
    \end{align}
    同様に$p_1$は
    \begin{align}
        p_1 = \dfrac{1 - \Re \bra{\psi}U\ket{\psi}}{2}
    \end{align}
\end{frame}

\end{document}
